\documentclass[10pt,letterpaper]{article}

\usepackage[T1]{fontenc}
\usepackage[utf8]{inputenc}
\usepackage[margin=0.82in,headheight=14pt]{geometry}
\usepackage{amsmath,amssymb,amsthm,mathtools}
\usepackage{newtxtext}
\usepackage{newtxmath}
\DeclareMathAlphabet{\mathcal}{OMS}{cmsy}{m}{n}
\usepackage{bm}
\usepackage{microtype}
\usepackage{booktabs}
\usepackage{array}
\usepackage{enumitem}
\usepackage{xcolor}
\usepackage[most]{tcolorbox}
\usepackage{fancyhdr}
\usepackage[colorlinks=true,linkcolor=blue!65!black,citecolor=blue!65!black,urlcolor=blue!65!black]{hyperref}

\emergencystretch=1.5em

\definecolor{navy}{HTML}{1B4F72}
\definecolor{crimson}{HTML}{922B21}
\definecolor{forest}{HTML}{145A32}
\definecolor{slate}{HTML}{34495E}
\definecolor{softblue}{HTML}{EBF5FB}
\definecolor{softgreen}{HTML}{EAFAF1}
\definecolor{softamber}{HTML}{FEF9E7}
\definecolor{amberborder}{HTML}{B7950B}

\newtcolorbox{targeteq}[1]{
  breakable,enhanced,lines before break=2,
  colback=softblue,colframe=navy,fonttitle=\bfseries\small,
  title={Exact Target in JOSA A Submission: #1},
  left=8pt,right=8pt,top=5pt,bottom=5pt,before skip=8pt,after skip=8pt
}

\newtcolorbox{pencheck}[1][Pen-and-Paper Worked Numerical Check]{
  breakable,enhanced,lines before break=2,
  colback=softgreen,colframe=forest,fonttitle=\bfseries\small,
  title={#1},
  left=8pt,right=8pt,top=5pt,bottom=5pt,before skip=8pt,after skip=8pt
}

\newtcolorbox{auditbox}[1][Sign, Convention \& Boundary Audit]{
  breakable,enhanced,lines before break=2,
  colback=softamber,colframe=amberborder,fonttitle=\bfseries\small,
  title={#1},
  left=8pt,right=8pt,top=5pt,bottom=5pt,before skip=8pt,after skip=8pt
}

\pagestyle{fancy}
\fancyhf{}
\fancyhead[L]{\footnotesize\color{slate}\textit{JOSA A Submission (V15 Paper 1) --- Complete Pen-and-Paper Workbook}}
\fancyhead[R]{\footnotesize\color{slate}\textit{S. Ayala}}
\fancyfoot[C]{\footnotesize\thepage}
\renewcommand{\headrulewidth}{0.4pt}

\newcommand{\braket}[2]{\langle #1 | #2 \rangle}
\newcommand{\ket}[1]{\lvert #1 \rangle}
\newcommand{\bra}[1]{\langle #1 \rvert}
\newcommand{\norm}[1]{\lVert #1 \rVert}
\newcommand{\abs}[1]{\lvert #1 \rvert}
\newcommand{\dd}{\mathrm{d}}
\newcommand{\josa}[1]{\textbf{[JOSA~A Eq.~(#1)]}}

\numberwithin{equation}{section}

\title{\vspace{-1.0cm}\textbf{\Large Complete Pen-and-Paper Mathematical Workbook for the JOSA A Submission (V15 Paper 1)}\\[0.35em]
\large\textit{Zero-Skipped-Steps Derivations of Every Numbered Equation (1)--(26), Inline Formula, and Table in\\
``Continuous Hilbert frame theory, paraxial propagation invariance, and exact Laguerre--Gaussian expansion of finite-width perfect optical vortex beams''}}
\author{Sebasti\'an Ayala\thanks{\texttt{ayala.phys.ds@gmail.com}; ORCID: \href{https://orcid.org/0009-0003-7348-3883}{0009-0003-7348-3883}}}
\date{Exact JOSA A Submission Edition --- November 24, 2026}

\begin{document}
\maketitle
\thispagestyle{plain}

\begin{abstract}
\noindent This workbook is indexed one-to-one to the exact numbered equations \textbf{Eqs.~(1)--(26)} (Sections~2--6 and Appendix~A) of the finalized \emph{JOSA~A} manuscript (\texttt{Manuscript\_JOSA\_A/main\_paper1\_josaa.pdf}). Every algebraic step, polar coordinate Jacobian, analytic continuation of Bessel integrals, Fourier normalization factor, Fr\'echet variation, Kummer factorial cancellation, WKB integral, Szeg\H{o} Toeplitz symbol, and hand-worked numerical check is written out in full so that the entire paper can be reproduced from scratch with pen and paper. All symbolic identities and numerical checks are independently verified in \texttt{Verification\_Scripts/verify\_paper1\_v15.py}.
\end{abstract}

\tableofcontents
\vspace{0.5em}

\section*{Exact 1-to-1 Cross-Reference Map to \texttt{main\_paper1\_josaa.pdf}}
\addcontentsline{toc}{section}{Exact 1-to-1 Cross-Reference Map to main\_paper1\_josaa.pdf}

\begin{table}[ht!]
\centering
\scriptsize
\renewcommand{\arraystretch}{1.15}
\setlength{\tabcolsep}{4pt}
\begin{tabular}{@{}llll@{}}
\toprule
\textbf{JOSA A Reference} & \textbf{Manuscript Page} & \textbf{Mathematical / Physical Content} & \textbf{Workbook Section} \\
\midrule
Eq.~(1), Eq.~(2) & p.~2 & Physical PVB $\Phi_{l,R}$, scaling $C_0$, dimensions, exact norm $\norm{\Phi_{l,R}}_{L^2}^2$ & Sec.~\ref{sec:s2_norm} \\
Eq.~(3) & p.~3 & Thin-ring Hankel limit, transcendental peak, $(l^2/8)(w/R)^4 R$ shift & Sec.~\ref{sec:s2_peak} \\
Eq.~(4) & p.~3 & $2f$ optical Fourier transform of $E_{\mathrm{BG}}$, $R, w, E_0$, two-zone propagation & Sec.~\ref{sec:s2_2f} \\
Eq.~(5), Eq.~(6) & pp.~3--4 & Sector decoupling, mutual overlap $\braket{\Phi_{l,R}}{\Phi_{l',R'}}$, universal $\gamma_l(R,R')$ & Sec.~\ref{sec:s3_overlap} \\
Eq.~(7), Eq.~(8), Eq.~(9) & p.~4 & 2D FT $\widetilde{\Phi}_{l,R}$, Hankel unitarity ($B=1$), convolution $S_w = G_w * \cdot$ & Sec.~\ref{sec:s3_frame} \\
Eq.~(10), Table~1 & pp.~4--5 & Sharp Paley--Wiener bounds $A = e^{-w^2 k_{\max}^2/2}, B=1$, annulus limits & Sec.~\ref{sec:s3_pw} \\
Eq.~(11), Eq.~(12) & p.~5 & Complex waist $w_c^2(z)$ and closed-form Fresnel PVB $\Phi_{l,R}^{(z)}(r,\phi)$ & Sec.~\ref{sec:s4_fresnel} \\
Eq.~(13), Eq.~(14), Eq.~(15) & pp.~5--6 & Commutation $[S_w,\mathcal{U}_z]=0$, $\sigma_m^{+1}$ analyzer, Gouy projection & Sec.~\ref{sec:s4_commute} \\
Eq.~(16), Eq.~(17) & pp.~6--7 & Tikhonov Fr\'echet derivative, $S_{w,\alpha_{\mathrm{reg}}}^{-1}$, AM--GM noise peak & Sec.~\ref{sec:s5_tikhonov} \\
Eq.~(18) & p.~7 & Parseval $\mathrm{Bias}^2 + \mathrm{Var}$ closed form, noisy $\mathrm{LG}_{2,1}$ benchmark & Sec.~\ref{sec:s5_parseval} \\
Eqs.~(19)--(20), (25)--(26) & pp.~8, 10 & Exact $(p+1)$-term $c_{l,p}(R)$ via GR 6.631.1 \& Kummer transformation & Sec.~\ref{sec:s6_clp} \\
Eq.~(21), Eq.~(22), Eq.~(23) & pp.~8--9 & Washout $\mathcal{F}$, exact $p=0$ root, rational $w=w_0$ checks, WKB arcsine $0.5724$ & Sec.~\ref{sec:s6_washout} \\
Eq.~(24), Table~2 & p.~9 & Unit-normalized $[\gamma]_{nm}$, Szeg\H{o} symbol $\vartheta_3(\theta/2,q)$, $\kappa(T_\infty) = \vartheta_3/\vartheta_4$ & Sec.~\ref{sec:s6_gram} \\
\bottomrule
\end{tabular}
\end{table}

% ==============================================================================
\section{JOSA A Section 2: Conventions, Exact Norm, Peak Shift, and $2f$ Bridge}
\label{sec:s2}

\subsection{Dimensions, Frame-vs.-State Scaling, and Exact $L^2(\mathbb{R}^2)$ Norm --- \josa{1}--\josa{2}}
\label{sec:s2_norm}

\begin{targeteq}{JOSA A Eqs.~(1) and (2) (p.~2)}
\[
  \Phi_{l,R}(r,\phi) = C_0 e^{il\phi}\exp\!\left(-\frac{r^2+R^2}{w^2}\right) I_{\abs{l}}\!\left(\frac{2Rr}{w^2}\right), \quad C_0 = \frac{1}{w^2}\sqrt{\frac{2}{\pi}}, \qquad
  \norm{\Phi_{l,R}}_{L^2}^2 = \frac{1}{w^2} e^{-R^2/w^2} I_{\abs{l}}\!\left(\frac{R^2}{w^2}\right) < \frac{1}{w^2}.
\]
\end{targeteq}

\paragraph*{Step 1: Complete dimensional bookkeeping.}
Denote the physical length dimension by $[\ell_0]$. Since $[r] = [R] = [w] = \ell_0$, the exponents $r^2/w^2$, $R^2/w^2$, and the Bessel argument $2Rr/w^2$ are all strictly dimensionless. Hence $I_{\abs{l}}(2Rr/w^2)$ is dimensionless, and
\begin{equation}
  [\Phi_{l,R}] = [C_0] = \left[\frac{1}{w^2}\sqrt{\frac{2}{\pi}}\right] = \ell_0^{-2}.
\end{equation}
For any unit-$L^2$-normalized physical wavefield $\psi \in L^2(\mathbb{R}^2)$, $\int_{\mathbb{R}^2}\abs{\psi(\mathbf{r})}^2\,\dd^2\mathbf{r} = 1$ with $[\dd^2\mathbf{r}] = [r\,\dd r\,\dd\phi] = \ell_0^2$, so $[\psi] = \ell_0^{-1}$. Therefore the analysis projection $(T\psi)(l,R) = \braket{\Phi_{l,R}}{\psi}$ has dimension
\begin{equation}
  [(T\psi)(l,R)] = [\Phi_{l,R}^*]\,[\psi]\,[\dd^2\mathbf{r}] = \ell_0^{-2}\cdot\ell_0^{-1}\cdot\ell_0^2 = \ell_0^{-1}.
\end{equation}
Integrating its squared modulus over the polar radial measure $\dd\mu(l,R) = R\,\dd R$ ($[\dd\mu] = \ell_0^2$) gives
\begin{equation}
  \left[\int_0^\infty R\,\dd R\,\abs{\braket{\Phi_{l,R}}{\psi}}^2\right] = \ell_0^2\cdot\ell_0^{-2} = 1\quad\text{(dimensionless)},
\end{equation}
matching $[\norm{\psi}_{L^2}^2] = 1$. Meanwhile the unit-normalized physical PVB state $\psi_{l,R} \equiv \Phi_{l,R}/\norm{\Phi_{l,R}}_{L^2}$ satisfies $[\norm{\Phi_{l,R}}_{L^2}] = ([\Phi_{l,R}]^2[\dd^2\mathbf{r}])^{1/2} = (\ell_0^{-4}\ell_0^2)^{1/2} = \ell_0^{-1}$, so $[\psi_{l,R}] = \ell_0^{-2}/\ell_0^{-1} = \ell_0^{-1}$.

\paragraph*{Step 2: Analytic continuation of Weber's second exponential integral.}
Gradshteyn--Ryzhik formula 6.633.2 states that for $\operatorname{Re}(p) > 0$, $\operatorname{Re}(\nu) > -1$, and $\alpha,\beta \in \mathbb{C}$,
\begin{equation}
  \int_0^\infty x\,e^{-p x^2} J_\nu(\alpha x) J_\nu(\beta x)\,\dd x = \frac{1}{2p}\exp\!\left(-\frac{\alpha^2 + \beta^2}{4p}\right) I_\nu\!\left(\frac{\alpha\beta}{2p}\right).
  \label{eq:wb_weber_J}
\end{equation}
To replace $J_\nu$ with $I_\nu$, recall the definition $I_\nu(z) = i^{-\nu} J_\nu(iz)$, i.e.\ $J_\nu(iz) = i^\nu I_\nu(z)$. Set $\alpha = ia$ and $\beta = ib$ ($a, b > 0$) with integer $\nu = \abs{l} \ge 0$:
\begin{align}
  \text{LHS} &= \int_0^\infty x\,e^{-px^2}\bigl[i^{\abs{l}} I_{\abs{l}}(ax)\bigr]\bigl[i^{\abs{l}} I_{\abs{l}}(bx)\bigr]\,\dd x = (-1)^{\abs{l}} \int_0^\infty x\,e^{-px^2} I_{\abs{l}}(ax) I_{\abs{l}}(bx)\,\dd x, \\
  \text{RHS} &= \frac{1}{2p}\exp\!\left(-\frac{-a^2 - b^2}{4p}\right) I_{\abs{l}}\!\left(-\frac{ab}{2p}\right) = (-1)^{\abs{l}}\frac{1}{2p}\exp\!\left(\frac{a^2+b^2}{4p}\right) I_{\abs{l}}\!\left(\frac{ab}{2p}\right),
\end{align}
where we used $I_{\abs{l}}(-y) = (-1)^{\abs{l}} I_{\abs{l}}(y)$ for integer $\abs{l}$. Dividing both sides by $(-1)^{\abs{l}}$ gives our master modified Weber identity:
\begin{equation}
  \int_0^\infty x\,e^{-p x^2} I_{\abs{l}}(ax) I_{\abs{l}}(bx)\,\dd x = \frac{1}{2p}\exp\!\left(\frac{a^2+b^2}{4p}\right) I_{\abs{l}}\!\left(\frac{ab}{2p}\right).
  \label{eq:wb_weber_I}
\end{equation}

\paragraph*{Step 3: Exact squared norm \josa{2}.}
Integrating $\abs{\Phi_{l,R}(r,\phi)}^2 = C_0^2 e^{-2R^2/w^2} e^{-2r^2/w^2}[I_{\abs{l}}(2Rr/w^2)]^2$ over $\phi \in [0,2\pi)$ gives $2\pi$:
\begin{equation}
  \norm{\Phi_{l,R}}_{L^2}^2 = 2\pi C_0^2 e^{-2R^2/w^2} \int_0^\infty r\,e^{-(2/w^2)r^2} I_{\abs{l}}\!\left(\frac{2R}{w^2}r\right) I_{\abs{l}}\!\left(\frac{2R}{w^2}r\right)\dd r.
\end{equation}
Identify $p = 2/w^2$ and $a = b = 2R/w^2$ in Eq.~\eqref{eq:wb_weber_I}:
\begin{equation}
  \frac{1}{2p} = \frac{w^2}{4}, \qquad \frac{a^2+b^2}{4p} = \frac{2\cdot(4R^2/w^4)}{8/w^2} = \frac{R^2}{w^2}, \qquad \frac{ab}{2p} = \frac{4R^2/w^4}{4/w^2} = \frac{R^2}{w^2}.
\end{equation}
Since $2\pi C_0^2 = 2\pi\cdot\frac{2}{\pi w^4} = \frac{4}{w^4}$, multiplying out cancels the $4$'s and combines $e^{-2R^2/w^2} e^{+R^2/w^2} = e^{-R^2/w^2}$:
\begin{equation}
  \norm{\Phi_{l,R}}_{L^2}^2 = \frac{4}{w^4} e^{-2R^2/w^2} \cdot \frac{w^2}{4} e^{+R^2/w^2} I_{\abs{l}}\!\left(\frac{R^2}{w^2}\right) = \frac{1}{w^2} e^{-R^2/w^2} I_{\abs{l}}\!\left(\frac{R^2}{w^2}\right).
\end{equation}
To prove the strict upper bound $< w^{-2}$ for all $R > 0$, use the integral representation $I_n(x) = \frac{1}{\pi}\int_0^\pi e^{x\cos\theta}\cos(n\theta)\,\dd\theta$:
\begin{equation}
  e^{-x} I_n(x) = \frac{1}{\pi}\int_0^\pi e^{-x(1-\cos\theta)}\cos(n\theta)\,\dd\theta \le \frac{1}{\pi}\int_0^\pi e^{-x(1-\cos\theta)}\,\dd\theta < \frac{1}{\pi}\int_0^\pi 1\,\dd\theta = 1\quad (x > 0),
\end{equation}
because $1 - \cos\theta > 0$ on $(0,\pi)$ so $e^{-x(1-\cos\theta)} < 1$ strictly on $(0,\pi)$! Thus $\norm{\Phi_{l,R}}_{L^2}^2 < w^{-2}$ for all $R > 0$ and all $l \in \mathbb{Z}$.

\begin{pencheck}[Hand Check at $R = 2\ell_0$, $w = 0.5\ell_0$, $l = 1$]
Here $x = R^2/w^2 = 4/0.25 = 16$, $w^{-2} = 4\ell_0^{-2}$. Since $e^{-16} I_1(16) = 0.097350$, Eq.~\josa{2} gives $\norm{\Phi_{1,2\ell_0}}_{L^2}^2 = 4\times 0.097350 = 0.3894\,\ell_0^{-2} < 4\,\ell_0^{-2}$.
\end{pencheck}

\subsection{Thin-Ring Asymptotics and Self-Consistent Differential Peak Shift --- \josa{3}}
\label{sec:s2_peak}

\begin{targeteq}{JOSA A Eq.~(3) and Inline Asymptotics (pp.~2--3)}
\[
  \Phi_{l,R}(r,\phi) \approx \frac{e^{il\phi}}{w\sqrt{2\pi^2 Rr}}\exp\!\left[-\frac{(r-R)^2}{w^2}\right], \qquad
  \Delta r_{\mathrm{peak}}(l) \equiv r_{\mathrm{peak}}(l) - r_{\mathrm{peak}}(0) \approx \frac{l^2}{8}\left(\frac{w}{R}\right)^4 R.
\]
\end{targeteq}

\paragraph*{Step 1: Thin-ring Hankel reduction.}
For $x = 2Rr/w^2 \gg 1$, Hankel's large-argument series (with $\mu \equiv 4l^2 - 1$) is
\begin{equation}
  I_{\abs{l}}(x) = \frac{e^x}{\sqrt{2\pi x}}\left[1 - \frac{\mu}{8x} + \frac{\mu(\mu-8)}{128 x^2} + \mathcal{O}(x^{-3})\right].
  \label{eq:wb_hankel}
\end{equation}
Inserting the leading term $e^x/\sqrt{2\pi x}$ into \josa{1} gives:
\begin{align}
  \exp\!\left(-\frac{r^2+R^2}{w^2}\right) e^{2Rr/w^2} &= \exp\!\left(-\frac{r^2 - 2Rr + R^2}{w^2}\right) = \exp\!\left[-\frac{(r-R)^2}{w^2}\right], \\
  C_0 \frac{1}{\sqrt{2\pi(2Rr/w^2)}} &= \frac{1}{w^2}\sqrt{\frac{2}{\pi}}\cdot\frac{w}{2\sqrt{\pi Rr}} = \frac{1}{w\sqrt{2\pi^2 Rr}}.
\end{align}
Near $r \to 0$, $I_{\abs{l}}(x) \sim (x/2)^{\abs{l}}/\abs{l}! \propto r^{\abs{l}}$, so the exact field vanishes regularly as $r^{\abs{l}}$ (for $l \neq 0$), whereas the thin-ring approximation diverges unphysically as $r^{-1/2}$.

\paragraph*{Step 2: Exact transcendental peak condition.}
The radial intensity $I(r) = \abs{\Phi_{l,R}(r)}^2 = C_0^2 e^{-2(r^2+R^2)/w^2}[I_{\abs{l}}(2Rr/w^2)]^2$ is strictly positive on $r > 0$. Taking its logarithm and differentiating with respect to $r$:
\begin{equation}
  \frac{\dd}{\dd r}\ln I(r) = -\frac{4r}{w^2} + 2\cdot\frac{2R}{w^2}\,\frac{I'_{\abs{l}}(x)}{I_{\abs{l}}(x)} = \frac{4}{w^2}\left[-r + R\frac{I'_{\abs{l}}(x)}{I_{\abs{l}}(x)}\right] = 0 \iff r_{\mathrm{peak}}(l) = R\frac{I'_{\abs{l}}(x_{\mathrm{peak}})}{I_{\abs{l}}(x_{\mathrm{peak}})},
\end{equation}
where $x_{\mathrm{peak}} = 2R\,r_{\mathrm{peak}}(l)/w^2$.

\paragraph*{Step 3: Self-consistent perturbation expansion to $\mathcal{O}((w/R)^4)$.}
Compute the logarithmic derivative of Eq.~\eqref{eq:wb_hankel} explicitly:
\begin{align}
  \ln I_{\abs{l}}(x) &= x - \frac{1}{2}\ln(2\pi) - \frac{1}{2}\ln x + \ln\!\left(1 - \frac{\mu}{8x} + \mathcal{O}(x^{-2})\right) = x - \frac{1}{2}\ln(2\pi) - \frac{1}{2}\ln x - \frac{\mu}{8x} + \mathcal{O}(x^{-2}), \\
  \frac{I'_{\abs{l}}(x)}{I_{\abs{l}}(x)} &= \frac{\dd}{\dd x}\ln I_{\abs{l}}(x) = 1 - \frac{1}{2x} + \frac{\mu}{8x^2} + \mathcal{O}(x^{-3}).
  \label{eq:wb_logderiv}
\end{align}
Let $\delta \equiv (w/R)^2 \ll 1$, and write $r_{\mathrm{peak}}(l)/R = 1 + c_1\delta + c_2\delta^2 + \mathcal{O}(\delta^3)$. Then
\begin{equation}
  x_{\mathrm{peak}}^{-1} = \frac{w^2}{2R\,r_{\mathrm{peak}}} = \frac{\delta}{2}\left(1 + c_1\delta + \mathcal{O}(\delta^2)\right)^{-1} = \frac{\delta}{2}\left(1 - c_1\delta + \mathcal{O}(\delta^2)\right) = \frac{\delta}{2} - \frac{c_1}{2}\delta^2 + \mathcal{O}(\delta^3),
\end{equation}
and $x_{\mathrm{peak}}^{-2} = \delta^2/4 + \mathcal{O}(\delta^3)$. Substitute both into Eq.~\eqref{eq:wb_logderiv}:
\begin{align}
  1 + c_1\delta + c_2\delta^2 &= 1 - \frac{1}{2}\left(\frac{\delta}{2} - \frac{c_1}{2}\delta^2\right) + \frac{\mu}{8}\left(\frac{\delta^2}{4}\right) + \mathcal{O}(\delta^3) \nonumber \\
  &= 1 - \frac{1}{4}\delta + \left(\frac{c_1}{4} + \frac{\mu}{32}\right)\delta^2 + \mathcal{O}(\delta^3).
\end{align}
Equating like powers of $\delta$ (notice how the self-consistent $-c_1\delta^2/2$ inside $-1/(2x_{\mathrm{peak}})$ shifts the quadratic coefficient!):
\begin{align}
  \mathcal{O}(\delta^1): &\quad c_1 = -\frac{1}{4}, \\
  \mathcal{O}(\delta^2): &\quad c_2 = \frac{-1/4}{4} + \frac{4l^2 - 1}{32} = -\frac{2}{32} + \frac{4l^2 - 1}{32} = \frac{4l^2 - 3}{32}.
\end{align}
Therefore,
\begin{equation}
  r_{\mathrm{peak}}(l) - R = -\frac{w^2}{4R} + \frac{4l^2 - 3}{32}\frac{w^4}{R^3} + \mathcal{O}\!\left(\frac{w^6}{R^5}\right), \qquad
  r_{\mathrm{peak}}(0) - R = -\frac{w^2}{4R} - \frac{3}{32}\frac{w^4}{R^3} + \mathcal{O}\!\left(\frac{w^6}{R^5}\right).
\end{equation}
Subtracting $r_{\mathrm{peak}}(0)$ from $r_{\mathrm{peak}}(l)$ cancels both $-\frac{w^2}{4R}$ and $-\frac{3}{32}\frac{w^4}{R^3}$, leaving \josa{3}:
\begin{equation}
  \Delta r_{\mathrm{peak}}(l) = r_{\mathrm{peak}}(l) - r_{\mathrm{peak}}(0) = \frac{4l^2}{32}\frac{w^4}{R^3} + \mathcal{O}\!\left(\frac{w^6}{R^5}\right) = \frac{l^2}{8}\left(\frac{w}{R}\right)^4 R.
\end{equation}

\begin{pencheck}[Hand Arithmetic at $w/R = 0.10$, $l = 4$ and $l = 6$]
At $w/R = 0.10$, $(w/R)^4 = 10^{-4}$: for $l = 4$, $(16/8)\times 10^{-4} = 2.000\times 10^{-4}$ (exact root: $2.024\times 10^{-4}$, $1.20\%$ diff.); for $l = 6$, $(36/8)\times 10^{-4} = 4.500\times 10^{-4}$ (exact root: $4.552\times 10^{-4}$, $1.14\%$ diff.).
\end{pencheck}

\subsection{Optical $2f$ Fourier Transform Bridge --- \josa{4}}
\label{sec:s2_2f}

\begin{targeteq}{JOSA A Eq.~(4) (p.~3)}
\[
  R = \frac{k_r f}{k_0}, \qquad w = \frac{2f}{k_0 w_g}, \qquad E_0 = \frac{1}{w w_g}\sqrt{\frac{2}{\pi}} = \frac{k_0}{f\sqrt{2\pi}}.
\]
\end{targeteq}

Place $E_{\mathrm{BG}}(\rho,\theta) = E_0 J_l(k_r\rho) e^{-\rho^2/w_g^2} e^{il\theta}$ at the front focal plane ($d_{\mathrm{in}} = f$) of a thin spherical lens of focal length $f$. By standard Fourier optics, the field at the back focal plane ($z = 0$) is the exact 2D optical Fourier transform:
\begin{equation}
  E_{\mathrm{out}}(r,\phi) = \frac{e^{2ik_0 f}}{i\lambda f}\int_0^\infty \rho\,\dd\rho \int_0^{2\pi} \dd\theta\, E_0 J_l(k_r\rho) e^{-\rho^2/w_g^2} e^{il\theta} \exp\!\left[-\frac{ik_0}{f} r\rho\cos(\theta-\phi)\right].
\end{equation}
\paragraph*{Step 1: Azimuthal integral via Jacobi--Anger.}
Let $\psi = \theta - \phi$ ($\theta = \phi + \psi$) and $\alpha = k_0 r\rho/f$. Using $e^{-i\alpha\cos\psi} = \sum_{n=-\infty}^\infty (-i)^n J_n(\alpha) e^{in\psi}$ and $\int_0^{2\pi} e^{i(l+n)\psi}\,\dd\psi = 2\pi\delta_{n,-l}$:
\begin{equation}
  \int_0^{2\pi} e^{il\theta} e^{-i\alpha\cos(\theta-\phi)}\,\dd\theta = e^{il\phi}\sum_{n=-\infty}^\infty (-i)^n J_n(\alpha)\,2\pi\delta_{n,-l} = 2\pi (-i)^{-l} J_{-l}(\alpha) e^{il\phi}.
\end{equation}
Check the sign identity for ALL integers $l \in \mathbb{Z}$ using $J_{-m}(\alpha) = (-1)^m J_m(\alpha)$ and $(-i)^{-m} = i^m = (-1)^m(-i)^m$:
\begin{itemize}[leftmargin=1.5em,itemsep=1pt,topsep=2pt]
  \item If $l = +m \ge 0$: $(-i)^{-m} J_{-m}(\alpha) = [(-1)^m(-i)^m][(-1)^m J_m(\alpha)] = (-i)^m J_m(\alpha) = (-i)^{\abs{l}} J_{\abs{l}}(\alpha)$.
  \item If $l = -m < 0$: $(-i)^{+m} J_{+m}(\alpha) = (-i)^{\abs{l}} J_{\abs{l}}(\alpha)$.
\end{itemize}
For signed $J_l(k_r\rho)$ in the input, note $J_{-m}(k_r\rho) = (-1)^m J_m(k_r\rho)$, and $(-1)^m (-i)^m = (+i)^m = (-i)^{-m} = (-i)^l$, so $(-i)^{\abs{l}} J_l(k_r\rho) J_{\abs{l}}(\alpha) = (-i)^l J_{\abs{l}}(k_r\rho) J_{\abs{l}}(\alpha)$ for all $l \in \mathbb{Z}$!
Substituting $\frac{2\pi}{i\lambda f} = \frac{-i k_0}{f}$ gives global phase $(-i)(-i)^l e^{2ik_0 f} = (-i)^{l+1} e^{2ik_0 f}$:
\begin{equation}
  E_{\mathrm{out}}(r,\phi) = (-i)^{l+1} e^{2ik_0 f} \frac{k_0 E_0}{f} e^{il\phi} \int_0^\infty \rho\,e^{-\rho^2/w_g^2} J_{\abs{l}}(k_r\rho) J_{\abs{l}}\!\left(\frac{k_0 r}{f}\rho\right)\dd\rho.
\end{equation}
\paragraph*{Step 2: Radial integral via Eq.~\eqref{eq:wb_weber_J}.}
Set $p = 1/w_g^2$, $\alpha = k_r$, $\beta = k_0 r/f$, $\nu = \abs{l}$ in Eq.~\eqref{eq:wb_weber_J}:
\begin{equation}
  \int_0^\infty \rho\,e^{-\rho^2/w_g^2} J_{\abs{l}}(k_r\rho) J_{\abs{l}}\!\left(\frac{k_0 r}{f}\rho\right)\dd\rho
  = \frac{w_g^2}{2}\exp\!\left[-\frac{w_g^2}{4}\left(k_r^2 + \frac{k_0^2 r^2}{f^2}\right)\right] I_{\abs{l}}\!\left(\frac{w_g^2 k_r k_0 r}{2f}\right).
\end{equation}
Matching exponents to $-(r^2+R^2)/w^2$ and $2Rr/w^2$:
\begin{align}
  \frac{w_g^2 k_0^2}{4f^2} = \frac{1}{w^2} &\implies w = \frac{2f}{k_0 w_g} \quad \left(w_g = \frac{2f}{k_0 w}\right), \\
  \frac{R^2}{w^2} = \frac{w_g^2 k_r^2}{4} = \frac{f^2 k_r^2}{k_0^2 w^2} &\implies R = \frac{k_r f}{k_0}, \qquad \frac{w_g^2 k_r k_0 r}{2f} = \frac{4f^2}{k_0^2 w^2}\cdot\frac{k_r k_0 r}{2f} = \frac{2(k_r f/k_0)r}{w^2} = \frac{2Rr}{w^2}.
\end{align}
Finally, equating the amplitude prefactor $\frac{k_0 E_0}{f}\cdot\frac{w_g^2}{2} = E_0 \frac{w_g}{w}$ to $C_0 = \frac{1}{w^2}\sqrt{\frac{2}{\pi}}$ yields \josa{4}:
\begin{equation}
  E_0 = \frac{1}{w w_g}\sqrt{\frac{2}{\pi}} = \frac{k_0}{2f}\sqrt{\frac{2}{\pi}} = \frac{k_0}{f\sqrt{2\pi}}.
\end{equation}
\paragraph*{Step 3: Two-zone ray vs.\ wave propagation.}
Before the lens, conical wavevectors propagate at angle $\sin\gamma = k_r/k_0 = R/f$; by right-triangle geometry, an on-axis circular block of radius $a_{\mathrm{obs}} \le R$ casts a conical shadow ending at $z_{\mathrm{heal}} = a_{\mathrm{obs}}/\tan\gamma \approx k_0 a_{\mathrm{obs}}/k_r$, while the finite illuminating waist $w_g$ limits the Bessel zone to $z_{\max} \approx k_0 w_g/k_r$. Past the focal plane ($z > 0$), the focused flat-phase ring of width $w$ has annular Rayleigh range $z_R^{\mathrm{ring}} = k_0 w^2/2$ (derived rigorously in Sec.~\ref{sec:s4_fresnel}).

% ==============================================================================
\section{JOSA A Section 3: Continuous Frame Operator and Sharp Paley--Wiener Bounds}
\label{sec:s3}

\subsection{Azimuthal--Radial Decoupling and Universal Radial Overlap --- \josa{5}--\josa{6}}
\label{sec:s3_overlap}

\begin{targeteq}{JOSA A Eqs.~(5) and (6) (pp.~3--4)}
\[
  \braket{\Phi_{l,R}}{\Phi_{l',R'}} = \delta_{l,l'}\frac{1}{w^2}e^{-(R^2+R'^2)/(2w^2)} I_{\abs{l}}\!\left(\frac{RR'}{w^2}\right), \qquad
  \gamma_l(R,R') = \frac{I_{\abs{l}}(RR'/w^2)}{\sqrt{I_{\abs{l}}(R^2/w^2)\,I_{\abs{l}}(R'^2/w^2)}}.
\]
\end{targeteq}

\paragraph*{Step 1: Exact overlap integral \josa{5}.}
\begin{equation}
  \braket{\Phi_{l,R}}{\Phi_{l',R'}} = C_0^2 e^{-(R^2+R'^2)/w^2}\underbrace{\int_0^{2\pi} e^{i(l'-l)\phi}\,\dd\phi}_{2\pi\delta_{l,l'}} \int_0^\infty r\,e^{-2r^2/w^2} I_{\abs{l}}\!\left(\frac{2Rr}{w^2}\right) I_{\abs{l'}}\!\left(\frac{2R'r}{w^2}\right)\dd r.
\end{equation}
Setting $l' = l$ and invoking Eq.~\eqref{eq:wb_weber_I} with $p = 2/w^2$, $a = 2R/w^2$, $b = 2R'/w^2$ ($\frac{1}{2p} = \frac{w^2}{4}$, $\frac{a^2+b^2}{4p} = \frac{R^2+R'^2}{2w^2}$, $\frac{ab}{2p} = \frac{RR'}{w^2}$) yields:
\begin{equation}
  \braket{\Phi_{l,R}}{\Phi_{l',R'}} = \delta_{l,l'}\frac{4}{w^4} e^{-(R^2+R'^2)/w^2}\cdot\frac{w^2}{4} e^{+(R^2+R'^2)/(2w^2)} I_{\abs{l}}\!\left(\frac{RR'}{w^2}\right) = \delta_{l,l'}\frac{1}{w^2} e^{-(R^2+R'^2)/(2w^2)} I_{\abs{l}}\!\left(\frac{RR'}{w^2}\right).
\end{equation}

\paragraph*{Step 2: Unit-normalized overlap \josa{6} and universal $l$-independent Gaussian asymptote.}
Dividing \josa{5} (at $l'=l$) by $\norm{\Phi_{l,R}}_{L^2}\norm{\Phi_{l,R'}}_{L^2} = w^{-2} e^{-(R^2+R'^2)/(2w^2)}\sqrt{I_{\abs{l}}(R^2/w^2)I_{\abs{l}}(R'^2/w^2)}$ cancels $w^{-2} e^{-(R^2+R'^2)/(2w^2)}$ completely, giving \josa{6}. Now substitute Eq.~\eqref{eq:wb_hankel} into numerator and denominator for $RR' \gg w^2$ ($\mu = 4l^2 - 1$):
\begin{align}
  \gamma_l(R,R') &= \frac{\frac{e^{RR'/w^2}}{\sqrt{2\pi RR'/w^2}}\left(1 - \frac{\mu w^2}{8RR'} + \mathcal{O}(w^4)\right)}{\left[\frac{e^{R^2/w^2}}{\sqrt{2\pi R^2/w^2}}\left(1 - \frac{\mu w^2}{8R^2}\right)\right]^{1/2}\left[\frac{e^{R'^2/w^2}}{\sqrt{2\pi R'^2/w^2}}\left(1 - \frac{\mu w^2}{8R'^2}\right)\right]^{1/2}}.
\end{align}
Notice the square-root prefactors cancel identically: $\frac{\sqrt{R}\sqrt{R'}}{\sqrt{RR'}} = 1$. The exponentials combine into $\exp(\frac{RR'}{w^2} - \frac{R^2+R'^2}{2w^2}) = \exp[-\frac{(R-R')^2}{2w^2}]$. Expanding $(1 - \epsilon)^{-1/2} = 1 + \frac{1}{2}\epsilon + \mathcal{O}(\epsilon^2)$:
\begin{align}
  \left(1 - \frac{\mu w^2}{8RR'}\right)\left(1 + \frac{\mu w^2}{16 R^2}\right)\left(1 + \frac{\mu w^2}{16 R'^2}\right) &= 1 + \frac{\mu w^2}{16}\left(\frac{1}{R^2} - \frac{2}{RR'} + \frac{1}{R'^2}\right) + \mathcal{O}(w^4) \nonumber \\
  &= 1 + \frac{(4l^2 - 1)w^2(R-R')^2}{16 R^2 R'^2} + \mathcal{O}(w^4).
\end{align}
Thus $\gamma_l(R,R') \approx e^{-(R-R')^2/(2w^2)}\{1 + [(4l^2-1)w^2(R-R')^2]/(16 R^2 R'^2)\}$, proving leading-order $l$-independence!

\subsection{2D Fourier Transform, Hankel Unitarity, and Convolution Kernel --- \josa{7}--\josa{9}}
\label{sec:s3_frame}

\begin{targeteq}{JOSA A Eqs.~(7), (8), and (9) (p.~4)}
\begin{align*}
  \widetilde{\Phi}_{l,R}(\mathbf{k}) &= \sqrt{2\pi}\,(-i)^{\abs{l}} e^{il\phi_k} e^{-w^2 k_\perp^2/4} J_{\abs{l}}(Rk_\perp), \\
  \norm{T\psi}_{L^2(\Omega,\mu)}^2 &= \int_{\mathbb{R}^2} e^{-w^2 k_\perp^2/2}\abs{\widetilde{\psi}(\mathbf{k})}^2\frac{\dd^2\mathbf{k}}{(2\pi)^2} \le \norm{\psi}_{L^2}^2, \qquad (S_w\psi)(\mathbf{r}) = (G_w * \psi)(\mathbf{r}).
\end{align*}
\end{targeteq}

\paragraph*{Step 1: Exact 2D Fourier transform \josa{7}.}
Using $\mathbf{k}\cdot\mathbf{r} = k_\perp r\cos(\phi - \phi_k)$ and the azimuthal Jacobi--Anger result $\int_0^{2\pi} e^{il\phi - ik_\perp r\cos(\phi-\phi_k)}\,\dd\phi = 2\pi(-i)^{\abs{l}} e^{il\phi_k} J_{\abs{l}}(k_\perp r)$ (valid for ALL $l \in \mathbb{Z}$, proven in Sec.~\ref{sec:s2_2f}):
\begin{equation}
  \widetilde{\Phi}_{l,R}(k_\perp,\phi_k) = 2\pi C_0 (-i)^{\abs{l}} e^{il\phi_k} e^{-R^2/w^2} \int_0^\infty r\,e^{-r^2/w^2} I_{\abs{l}}\!\left(\frac{2Rr}{w^2}\right) J_{\abs{l}}(k_\perp r)\,\dd r.
\end{equation}
In Eq.~\eqref{eq:wb_weber_I}, set $p = 1/w^2$, $a = 2R/w^2$, and analytically continue $b \to -i k_\perp$ using $I_{\abs{l}}(-i z) = (-i)^{\abs{l}} J_{\abs{l}}(z)$ on both sides:
\begin{equation}
  (-i)^{\abs{l}}\int_0^\infty r\,e^{-r^2/w^2} I_{\abs{l}}\!\left(\frac{2Rr}{w^2}\right) J_{\abs{l}}(k_\perp r)\,\dd r = \frac{w^2}{2}\exp\!\left(\frac{4R^2/w^4 - k_\perp^2}{4/w^2}\right)\cdot (-i)^{\abs{l}} J_{\abs{l}}(Rk_\perp).
\end{equation}
Cancel $(-i)^{\abs{l}}$ on both sides, observe $\frac{4R^2/w^4 - k_\perp^2}{4/w^2} = \frac{R^2}{w^2} - \frac{w^2 k_\perp^2}{4}$ so $e^{-R^2/w^2} e^{+R^2/w^2} = 1$, and evaluate $2\pi C_0(w^2/2) = \pi w^2\cdot w^{-2}\sqrt{2/\pi} = \sqrt{2\pi}$, yielding \josa{7}:
\begin{equation}
  \widetilde{\Phi}_{l,R}(\mathbf{k}) = \sqrt{2\pi}\,(-i)^{\abs{l}} e^{il\phi_k}\exp\!\left(-\frac{w^2 k_\perp^2}{4}\right) J_{\abs{l}}(R k_\perp).
\end{equation}

\paragraph*{Step 2: Hankel unitarity and Bessel bound $B = 1$ \josa{8}.}
Expand $\widetilde{\psi}(k_\perp,\phi_k) = \sum_{l'\in\mathbb{Z}}\widetilde{\psi}_{l'}(k_\perp) e^{il'\phi_k}$. Note that its $L^2$ norm satisfies
\begin{equation}
  \norm{\psi}_{L^2}^2 = \frac{1}{(2\pi)^2}\int_0^\infty k_\perp\,\dd k_\perp\int_0^{2\pi}\dd\phi_k\,\abs{\widetilde{\psi}(k_\perp,\phi_k)}^2 = \frac{1}{2\pi}\sum_{l\in\mathbb{Z}}\int_0^\infty k_\perp\,\dd k_\perp\,\abs{\widetilde{\psi}_l(k_\perp)}^2.
  \label{eq:wb_parseval_polar}
\end{equation}
By Plancherel's identity, $\braket{\Phi_{l,R}}{\psi} = (2\pi)^{-2}\int_0^\infty k_\perp\,\dd k_\perp\int_0^{2\pi}\dd\phi_k\,\widetilde{\Phi}_{l,R}^*\widetilde{\psi}$. Since $[\sqrt{2\pi}(-i)^{\abs{l}} e^{il\phi_k}]^* = \sqrt{2\pi}\,i^{\abs{l}} e^{-il\phi_k}$ and $\int_0^{2\pi} e^{i(l'-l)\phi_k}\,\dd\phi_k = 2\pi\delta_{l,l'}$, we obtain
\begin{equation}
  (T\psi)(l,R) = \frac{2\pi\sqrt{2\pi}}{(2\pi)^2}\,i^{\abs{l}}\int_0^\infty k_\perp\,\dd k_\perp\,e^{-w^2 k_\perp^2/4}\widetilde{\psi}_l(k_\perp) J_{\abs{l}}(Rk_\perp) = \frac{i^{\abs{l}}}{\sqrt{2\pi}}\,\mathrm{H}_{\abs{l}}\!\left[e^{-w^2 k_\perp^2/4}\widetilde{\psi}_l\right](R).
\end{equation}
By the Hankel completeness relation $\int_0^\infty R\,J_{\abs{l}}(Rk) J_{\abs{l}}(Rk')\,\dd R = k^{-1}\delta(k-k')$, the Hankel transform $\mathrm{H}_{\abs{l}}$ is an isometry: $\int_0^\infty R\,\dd R\,\abs{\mathrm{H}_{\abs{l}}[f](R)}^2 = \int_0^\infty k_\perp\,\dd k_\perp\,\abs{f(k_\perp)}^2$. Therefore:
\begin{equation}
  \norm{T\psi}_{L^2(\Omega,\mu)}^2 = \sum_{l\in\mathbb{Z}}\int_0^\infty R\,\dd R\,\abs{(T\psi)(l,R)}^2 = \frac{1}{2\pi}\sum_{l\in\mathbb{Z}}\int_0^\infty k_\perp\,\dd k_\perp\,e^{-w^2 k_\perp^2/2}\abs{\widetilde{\psi}_l(k_\perp)}^2 = \int_{\mathbb{R}^2} e^{-w^2 k_\perp^2/2}\abs{\widetilde{\psi}(\mathbf{k})}^2\frac{\dd^2\mathbf{k}}{(2\pi)^2}.
\end{equation}
Since $0 < e^{-w^2 k_\perp^2/2} \le 1$ for all $k_\perp \ge 0$, Eq.~\eqref{eq:wb_parseval_polar} immediately yields $\norm{T\psi}^2 \le \norm{\psi}_{L^2}^2$, proving \josa{8} with $B = 1$.

\paragraph*{Step 3: Polarization and Gaussian convolution kernel \josa{9}.}
For any bounded linear operator $T$, the polarization identity $\braket{T\phi}{T\psi} = \frac{1}{4}\sum_{m=0}^3 i^m \norm{T(\psi + i^m\phi)}^2$ applied to \josa{8} yields
\begin{equation}
  \braket{\phi}{S_w\psi} = \braket{T\phi}{T\psi}_{L^2(\Omega,\mu)} = \int_{\mathbb{R}^2} \widetilde{\phi}^*(\mathbf{k})\left[e^{-w^2 k_\perp^2/2}\widetilde{\psi}(\mathbf{k})\right]\frac{\dd^2\mathbf{k}}{(2\pi)^2}.
\end{equation}
By the 2D convolution theorem $\mathcal{F}\{G_w * \psi\}(\mathbf{k}) = \widetilde{G}_w(\mathbf{k})\widetilde{\psi}(\mathbf{k})$, compute $G_w(\mathbf{r}) = \mathcal{F}^{-1}\{e^{-w^2(k_x^2+k_y^2)/2}\}$ explicitly by completing the square in each 1D Gaussian integral:
\begin{equation}
  \int_{-\infty}^\infty e^{-\frac{w^2}{2}k_x^2 + i k_x x}\,\frac{\dd k_x}{2\pi} = \frac{e^{-x^2/(2w^2)}}{2\pi}\int_{-\infty}^\infty e^{-\frac{w^2}{2}(k_x - ix/w^2)^2}\,\dd k_x = \frac{e^{-x^2/(2w^2)}}{2\pi}\sqrt{\frac{2\pi}{w^2}} = \frac{1}{\sqrt{2\pi w^2}}e^{-x^2/(2w^2)}.
\end{equation}
Multiplying the $x$- and $y$-factors produces $G_w(\mathbf{r}) = \frac{1}{2\pi w^2}e^{-(x^2+y^2)/(2w^2)}$, proving \josa{9}. Because $S_w$ is unitarily equivalent via $\mathcal{F}$ to multiplication by the continuous function $\widetilde{G}_w(k_\perp) = e^{-w^2 k_\perp^2/2} \in (0,1]$, it is positive, self-adjoint, bounded with $\norm{S_w} = \sup\widetilde{G}_w = 1$, and has purely continuous spectrum $\sigma(S_w) = \overline{(0,1]} = [0,1]$ (hence non-compact on unconstrained $\mathbb{R}^2$, since compact self-adjoint operators have at most countable spectrum).

\subsection{Sharp Paley--Wiener Bounds and Fourier-Annulus Limits --- \josa{10}}
\label{sec:s3_pw}

\begin{targeteq}{JOSA A Eq.~(10) and Inline Sharpness Limits (p.~4)}
\[
  A(w,k_{\max}) = \exp\!\left(-\frac{w^2 k_{\max}^2}{2}\right), \qquad B = 1, \qquad \kappa(S_{w,k_{\max}}) = \exp\!\left(\frac{w^2 k_{\max}^2}{2}\right).
\]
\end{targeteq}

\paragraph*{Step 1: Unboundedness on $L^2(\mathbb{R}^2)$ vs.\ boundedness on $\mathcal{H}_{k_{\max}}$.}
In Fourier space, $(S_w^{-1}\Phi_{l,R})\,\widetilde{}\;(\mathbf{k}) = e^{+w^2 k_\perp^2/2}\widetilde{\Phi}_{l,R}(\mathbf{k}) = \sqrt{2\pi}(-i)^{\abs{l}} e^{il\phi_k} e^{+w^2 k_\perp^2/4} J_{\abs{l}}(R k_\perp)$. Since $J_{\abs{l}}(Rk_\perp) \sim \sqrt{2/(\pi R k_\perp)}\cos(Rk_\perp - \dots)$, its squared modulus grows exponentially as $e^{+w^2 k_\perp^2/2}/k_\perp$, so $S_w^{-1}\Phi_{l,R} \notin L^2(\mathbb{R}^2)$. Restricted to $\mathcal{H}_{k_{\max}} = \{\psi : \operatorname{supp}(\widetilde{\psi}) \subseteq \{k_\perp \le k_{\max}\}\}$, however, $e^{-w^2 k_{\max}^2/2} \le e^{-w^2 k_\perp^2/2} \le 1$ on $k_\perp \in [0, k_{\max}]$, so $A \ge e^{-w^2 k_{\max}^2/2}$ and $B \le 1$, and the inverse multiplier $e^{+w^2 k_\perp^2/2}\mathbf{1}_{\{k_\perp \le k_{\max}\}}$ is bounded by $A^{-1} = e^{+w^2 k_{\max}^2/2}$.

\paragraph*{Step 2: Explicit proof of sharpness.}
For $\epsilon \in (0, k_{\max})$, construct two unit-$L^2$-normalized Fourier-annulus wavefields in $\mathcal{H}_{k_{\max}}$:
\begin{equation}
  \widetilde{\psi}_\epsilon^{(+)}(\mathbf{k}) = \frac{2\pi}{\sqrt{\pi\epsilon^2}}\mathbf{1}_{\{0 < k_\perp \le \epsilon\}}, \qquad
  \widetilde{\psi}_\epsilon^{(-)}(\mathbf{k}) = \frac{2\pi}{\sqrt{\pi\epsilon(2k_{\max}-\epsilon)}}\mathbf{1}_{\{k_{\max}-\epsilon \le k_\perp \le k_{\max}\}}.
\end{equation}
Verify unit normalization of $\psi_\epsilon^{(-)}$: the annulus area is $\pi[k_{\max}^2 - (k_{\max}-\epsilon)^2] = \pi(2k_{\max}\epsilon - \epsilon^2) = \pi\epsilon(2k_{\max}-\epsilon)$, so
\begin{equation}
  \norm{\psi_\epsilon^{(-)}}_{L^2}^2 = \frac{1}{(2\pi)^2}\cdot\frac{4\pi^2}{\pi\epsilon(2k_{\max}-\epsilon)}\cdot\pi\epsilon(2k_{\max}-\epsilon) = 1.
\end{equation}
Now evaluate its exact Rayleigh quotient by elementary substitution $t = w^2 k_\perp^2/2$ ($\dd t = w^2 k_\perp\,\dd k_\perp$):
\begin{align}
  \braket{\psi_\epsilon^{(-)}}{S_w\psi_\epsilon^{(-)}} &= \frac{1}{(2\pi)^2}\cdot\frac{4\pi^2}{\pi\epsilon(2k_{\max}-\epsilon)}\int_{k_{\max}-\epsilon}^{k_{\max}} 2\pi k_\perp\,e^{-w^2 k_\perp^2/2}\,\dd k_\perp \nonumber \\
  &= \frac{2}{\epsilon(2k_{\max}-\epsilon)}\cdot\frac{1}{w^2}\left[-e^{-w^2 k_\perp^2/2}\right]_{k_{\max}-\epsilon}^{k_{\max}} = \frac{2\left[e^{-w^2(k_{\max}-\epsilon)^2/2} - e^{-w^2 k_{\max}^2/2}\right]}{w^2\epsilon(2k_{\max}-\epsilon)}.
\end{align}
Taking $\epsilon \to 0^+$ via L'H\^opital's rule (differentiating numerator and denominator with respect to $\epsilon$):
\begin{equation}
  \lim_{\epsilon\to 0^+} \frac{2\cdot w^2(k_{\max}-\epsilon)\,e^{-w^2(k_{\max}-\epsilon)^2/2}}{w^2(2k_{\max} - 2\epsilon)} = \frac{2k_{\max}}{2k_{\max}} e^{-w^2 k_{\max}^2/2} = e^{-w^2 k_{\max}^2/2} = A(w, k_{\max}).
\end{equation}
Similarly, $\braket{\psi_\epsilon^{(+)}}{S_w\psi_\epsilon^{(+)}} = \frac{2}{\epsilon^2 w^2}(1 - e^{-w^2\epsilon^2/2}) \to \frac{2}{\epsilon^2 w^2}\cdot\frac{w^2\epsilon^2}{2} = 1 = B$. Both bounds in \josa{10} are therefore strictly sharp.

% ==============================================================================
\section{JOSA A Section 4: Exact Paraxial Fresnel Propagation Invariance}
\label{sec:s4}

\subsection{Complex-Waist Closed-Form Propagated PVB --- \josa{11}--\josa{12}}
\label{sec:s4_fresnel}

\begin{targeteq}{JOSA A Eqs.~(11) and (12) (p.~5)}
\[
  w_c^2(z) = w^2\!\left(1 + i\frac{z}{z_R^{\mathrm{ring}}}\right), \quad z_R^{\mathrm{ring}} = \frac{k_0 w^2}{2}, \qquad
  \Phi_{l,R}^{(z)}(r,\phi) = \frac{1}{w_c^2(z)}\sqrt{\frac{2}{\pi}}\,e^{il\phi}\exp\!\left(-\frac{r^2+R^2}{w_c^2(z)}\right) I_{\abs{l}}\!\left(\frac{2Rr}{w_c^2(z)}\right).
\]
\end{targeteq}

\paragraph*{Step 1: Fresnel propagator in Fourier space.}
Taking the 2D transverse Fourier transform of the paraxial Helmholtz equation $(2ik_0\partial_z + \nabla_\perp^2)\psi = 0$ replaces $\nabla_\perp^2 \mapsto -k_\perp^2$:
\begin{equation}
  2ik_0\frac{\partial\widetilde{\psi}}{\partial z} - k_\perp^2\widetilde{\psi} = 0 \implies \frac{\partial\widetilde{\psi}}{\partial z} = -\frac{i k_\perp^2}{2k_0}\widetilde{\psi} \implies \widetilde{\psi}(\mathbf{k}, z) = e^{-iz k_\perp^2/(2k_0)}\widetilde{\psi}(\mathbf{k}, 0).
\end{equation}
Multiplying Eq.~\josa{7} by $e^{-iz k_\perp^2/(2k_0)}$ merges both $k_\perp^2$ exponents cleanly into \josa{11}:
\begin{equation}
  -\frac{w^2 k_\perp^2}{4} - \frac{i z k_\perp^2}{2k_0} = -\frac{k_\perp^2}{4}\left(w^2 + \frac{2iz}{k_0}\right) = -\frac{k_\perp^2}{4} w^2\!\left(1 + i\frac{z}{k_0 w^2/2}\right) = -\frac{w_c^2(z) k_\perp^2}{4}.
\end{equation}

\paragraph*{Step 2: Exact inverse 2D Fourier transform \josa{12}.}
Apply $\mathcal{F}^{-1}$ to $\widetilde{\Phi}_{l,R}^{(z)}(\mathbf{k}) = \sqrt{2\pi}(-i)^{\abs{l}} e^{il\phi_k}\allowbreak\exp(-w_c^2(z)k_\perp^2/4) J_{\abs{l}}(R k_\perp)$:
\begin{equation}
  \Phi_{l,R}^{(z)}(r,\phi) = \frac{\sqrt{2\pi}(-i)^{\abs{l}}}{(2\pi)^2}\int_0^\infty k_\perp\,\dd k_\perp\, e^{-w_c^2(z)k_\perp^2/4} J_{\abs{l}}(Rk_\perp)\int_0^{2\pi} \dd\phi_k\, e^{il\phi_k + i k_\perp r\cos(\phi_k-\phi)}.
\end{equation}
Evaluating the azimuthal integral via $e^{+i\alpha\cos\psi} = \sum_n i^n J_n(\alpha) e^{in\psi}$ gives $2\pi\,i^{\abs{l}} e^{il\phi} J_{\abs{l}}(r k_\perp)$ (matching the sign check of Sec.~\ref{sec:s2_2f} with $+i$ in place of $-i$). Since $(-i)^{\abs{l}} i^{\abs{l}} = 1$ and $\frac{2\pi\sqrt{2\pi}}{(2\pi)^2} = \frac{1}{\sqrt{2\pi}}$:
\begin{equation}
  \Phi_{l,R}^{(z)}(r,\phi) = \frac{e^{il\phi}}{\sqrt{2\pi}}\int_0^\infty k_\perp\,e^{-w_c^2(z)k_\perp^2/4} J_{\abs{l}}(Rk_\perp) J_{\abs{l}}(rk_\perp)\,\dd k_\perp.
\end{equation}
Because $\operatorname{Re}[w_c^2(z)/4] = w^2/4 > 0$ strictly for EVERY $z \in (-\infty,\infty)$, Weber's integral~\eqref{eq:wb_weber_J} applies unconditionally with complex $p = w_c^2(z)/4$, $\alpha = R$, $\beta = r$:
\begin{equation}
  \frac{1}{2p} = \frac{2}{w_c^2(z)}, \qquad -\frac{\alpha^2+\beta^2}{4p} = -\frac{r^2+R^2}{w_c^2(z)}, \qquad \frac{\alpha\beta}{2p} = \frac{2Rr}{w_c^2(z)}.
\end{equation}
Multiplying by $\frac{e^{il\phi}}{\sqrt{2\pi}}$ gives $\frac{2}{\sqrt{2\pi}} = \sqrt{\frac{2}{\pi}}$, yielding \josa{12} identically!

\paragraph*{Step 3: Real and imaginary parts of $1/w_c^2(z)$.}
Rationalizing the complex denominator $w_c^2(z) = w^2(1 + i z/z_R^{\mathrm{ring}})$ gives
\begin{equation}
  \frac{1}{w_c^2(z)} = \frac{1 - i z/z_R^{\mathrm{ring}}}{w^2\bigl[1 + (z/z_R^{\mathrm{ring}})^2\bigr]} = \frac{1}{w^2(z)} - \frac{i k_0}{2 R_c^{\mathrm{ring}}(z)},
\end{equation}
since $\frac{z/z_R^{\mathrm{ring}}}{w^2[1+(z/z_R^{\mathrm{ring}})^2]} = \frac{2z/(k_0 w^2)}{w^2(z^2+(z_R^{\mathrm{ring}})^2)/(z_R^{\mathrm{ring}})^2} = \frac{k_0 z}{2[z^2 + (z_R^{\mathrm{ring}})^2]} = \frac{k_0}{2 z[1+(z_R^{\mathrm{ring}}/z)^2]} = \frac{k_0}{2 R_c^{\mathrm{ring}}(z)}$!

\subsection{Commutation $[S_w,\mathcal{U}_z] = 0$ and Corrected Gouy Analyzer --- \josa{13}--\josa{15}}
\label{sec:s4_commute}

\begin{targeteq}{JOSA A Eqs.~(13), (14), and (15) (pp.~5--6)}
\begin{align*}
  [S_w, \mathcal{U}_z] = 0, \quad [P_{k_{\max}}, \mathcal{U}_z] = 0 &\implies S_w^{(z)} = \mathcal{U}_z S_w \mathcal{U}_z^\dagger = S_w, \\
  \phi_{l,R}^{(m)}(r,\phi) = \sigma_m e^{+ik_0 r^2/[2R_c(z_m)]}\Phi_{l,\sigma_m R}(r,\phi;\,\sigma_m w) &\implies \braket{\phi_{l,R}^{(m)}}{\psi_{p,l}(z_m)} = e^{-i(2p+\abs{l}+1)\vartheta_m} c_{l,p}(R).
\end{align*}
\end{targeteq}

\paragraph*{Step 1: Unitary invariance \josa{13}.}
In Fourier space, $(S_w\mathcal{U}_z\psi)\,\widetilde{}\;(\mathbf{k}) = e^{-w^2 k_\perp^2/2} e^{-iz k_\perp^2/(2k_0)}\widetilde{\psi}(\mathbf{k}) = (\mathcal{U}_z S_w\psi)\,\widetilde{}\;(\mathbf{k})$, and similarly for $P_{k_{\max}}$. Because $\abs{e^{-iz k_\perp^2/(2k_0)}} = 1$, $\mathcal{U}_z$ is unitary ($\mathcal{U}_z^\dagger = \mathcal{U}_{-z} = \mathcal{U}_z^{-1}$). Therefore the frame operator at plane $z$ is
\begin{equation}
  S_w^{(z)}\psi = \sum_{l\in\mathbb{Z}}\int_0^\infty R\,\dd R\,\braket{\mathcal{U}_z\Phi_{l,R}}{\psi}\,\mathcal{U}_z\Phi_{l,R} = \mathcal{U}_z\left(\sum_{l\in\mathbb{Z}}\int_0^\infty R\,\dd R\,\braket{\Phi_{l,R}}{\mathcal{U}_z^\dagger\psi}\Phi_{l,R}\right) = \mathcal{U}_z S_w \mathcal{U}_z^\dagger\psi = S_w\mathcal{U}_z\mathcal{U}_z^\dagger\psi = S_w\psi.
\end{equation}
Likewise, $\braket{\Phi_{l,R}^{(z)}}{\psi_{p,l}(z)} = \braket{\mathcal{U}_z\Phi_{l,R}}{\mathcal{U}_z\psi_{p,l}(0)} = \braket{\Phi_{l,R}}{\mathcal{U}_z^\dagger\mathcal{U}_z\psi_{p,l}(0)} = c_{l,p}(R)$.

\paragraph*{Step 2: Exact $\sigma_m^{+1}$ dimensional scaling \josa{14}--\josa{15}.}
Under Fresnel evolution $e^{-iz_m k_\perp^2/(2k_0)}$, a unit-norm $\mathrm{LG}_{p,l}$ beam of waist $w_0$ and Rayleigh range $z_R = k_0 w_0^2/2$ takes the exact form
\begin{equation}
  \psi_{p,l}(r,\phi,z_m) = \frac{1}{\sigma_m}\psi_{p,l}\!\left(\frac{r}{\sigma_m},\phi,0\right)\exp\!\left[+\frac{ik_0 r^2}{2 R_c(z_m)} - i(2p+\abs{l}+1)\vartheta_m\right],
\end{equation}
with $\sigma_m = \sqrt{1+(z_m/z_R)^2}$, $R_c(z_m) = z_m[1+(z_R/z_m)^2]$, and $\vartheta_m = \arctan(z_m/z_R)$. Now substitute $(R,w) \mapsto (\sigma_m R, \sigma_m w)$ inside \josa{1}:
\begin{equation}
  \Phi_{l,\sigma_m R}(r,\phi;\,\sigma_m w) = \frac{\sqrt{2/\pi}}{\sigma_m^2 w^2} e^{il\phi}\exp\!\left(-\frac{(r/\sigma_m)^2 + R^2}{w^2}\right) I_{\abs{l}}\!\left(\frac{2R(r/\sigma_m)}{w^2}\right) = \sigma_m^{-2}\,\Phi_{l,R}\!\left(\frac{r}{\sigma_m},\phi;\,w\right).
\end{equation}
Multiplying by $\sigma_m^{+1} e^{+ik_0 r^2/[2R_c(z_m)]}$ defines \josa{14}, which equals $\sigma_m^{-1} e^{+ik_0 r^2/[2R_c(z_m)]}\Phi_{l,R}(r/\sigma_m,\phi;\,w)$. Substituting $\mathbf{r} = \sigma_m\mathbf{r}'$ ($\dd^2\mathbf{r} = \sigma_m^2\,\dd^2\mathbf{r}'$) cancels the quadratic wavefront phase $e^{-ik_0 r^2/(2R_c)} e^{+ik_0 r^2/(2R_c)} = 1$ and the scaling factor $\sigma_m^{-1}\cdot\sigma_m^{-1}\cdot\sigma_m^2 = 1$, proving \josa{15}:
\begin{equation}
  \braket{\phi_{l,R}^{(m)}}{\psi_{p,l}(z_m)} = e^{-i(2p+\abs{l}+1)\vartheta_m}\int_{\mathbb{R}^2}\Phi_{l,R}^*(\mathbf{r}';w)\psi_{p,l}(\mathbf{r}',0)\,\dd^2\mathbf{r}' = e^{-i(2p+\abs{l}+1)\vartheta_m}\,c_{l,p}(R).
\end{equation}

\begin{auditbox}[Machine-Precision Verification of Eqs.~(11)--(15)]
In \texttt{Verification\_Scripts/verify\_paper1\_v15.py} (Suite~3): Eq.~\josa{12} agrees with 40-digit inverse Hankel quadrature to $5.27\times 10^{-23}$ and with $512\times 512$ 2D-FFT propagation to $4.73\times 10^{-16}$; unitary invariance holds to $6.27\times 10^{-16}$; and Eq.~\josa{15} holds to $5.78\times 10^{-16}$.
\end{auditbox}

% ==============================================================================
\section{JOSA A Section 5: Tikhonov Dual and Worked Noisy $\mathrm{LG}_{2,1}$ Deconvolution}
\label{sec:s5}

\subsection{Fr\'echet Derivative and AM--GM Noise-Gain Peak --- \josa{16}--\josa{17}}
\label{sec:s5_tikhonov}

\begin{targeteq}{JOSA A Eqs.~(16) and (17) (pp.~6--7)}
\[
  S_{w,\alpha_{\mathrm{reg}}}^{-1} = (S_w^2 + \alpha_{\mathrm{reg}}\mathbb{I})^{-1}S_w, \qquad
  \widetilde{H}_{w,\alpha_{\mathrm{reg}}}(k_\perp) = \frac{e^{-w^2 k_\perp^2/2}}{e^{-w^2 k_\perp^2} + \alpha_{\mathrm{reg}}}, \qquad
  \widetilde{H}_{\max} = \frac{1}{2\sqrt{\alpha_{\mathrm{reg}}}}\text{ at } k_{\mathrm{peak}} = \frac{1}{w}\sqrt{\ln\!\left(\frac{1}{\alpha_{\mathrm{reg}}}\right)}.
\]
\end{targeteq}

\paragraph*{Step 1: Variational Fr\'echet derivative \josa{16}.}
Expand $J[\psi + \epsilon\eta] = \norm{S_w(\psi + \epsilon\eta) - \psi_{\mathrm{meas}}}_{L^2}^2 + \alpha_{\mathrm{reg}}\norm{\psi + \epsilon\eta}_{L^2}^2$ in real $\epsilon$ using $S_w^\dagger = S_w$:
\begin{align}
  J[\psi + \epsilon\eta] &= J[\psi] + 2\epsilon\operatorname{Re}\bigl[\braket{S_w\eta}{S_w\psi - \psi_{\mathrm{meas}}} + \alpha_{\mathrm{reg}}\braket{\eta}{\psi}\bigr] + \epsilon^2\bigl(\norm{S_w\eta}^2 + \alpha_{\mathrm{reg}}\norm{\eta}^2\bigr) \nonumber \\
  &= J[\psi] + 2\epsilon\operatorname{Re}\braket{\eta}{(S_w^2 + \alpha_{\mathrm{reg}}\mathbb{I})\psi - S_w\psi_{\mathrm{meas}}} + \mathcal{O}(\epsilon^2).
\end{align}
Because the $\epsilon^2$ term is strictly positive for $\alpha_{\mathrm{reg}} > 0$ and $\eta \neq 0$, $J$ is strictly convex and its unique global minimizer satisfies $(S_w^2 + \alpha_{\mathrm{reg}}\mathbb{I})\hat{\psi}_{\alpha_{\mathrm{reg}}} = S_w\psi_{\mathrm{meas}}$, proving \josa{16}.

\paragraph*{Step 2: Spectral peak \josa{17} via AM--GM and calculus.}
Setting $g = \widetilde{G}_w(k_\perp) = e^{-w^2 k_\perp^2/2} \in (0,1]$ yields $\widetilde{H}_{w,\alpha_{\mathrm{reg}}}(g) = \frac{g}{g^2+\alpha_{\mathrm{reg}}} = (g + \alpha_{\mathrm{reg}}/g)^{-1}$. Differentiating with respect to $g$:
\begin{equation}
  \frac{\dd}{\dd g}\left(\frac{g}{g^2+\alpha_{\mathrm{reg}}}\right) = \frac{(g^2+\alpha_{\mathrm{reg}}) - 2g^2}{(g^2+\alpha_{\mathrm{reg}})^2} = \frac{\alpha_{\mathrm{reg}} - g^2}{(g^2+\alpha_{\mathrm{reg}})^2} = 0 \iff g_* = \sqrt{\alpha_{\mathrm{reg}}}.
\end{equation}
For $\alpha_{\mathrm{reg}} \in (0,1)$, $g_* = \sqrt{\alpha_{\mathrm{reg}}} \in (0,1)$ lies inside the range of $g(k_\perp)$, giving $\widetilde{H}_{\max} = \frac{\sqrt{\alpha_{\mathrm{reg}}}}{\alpha_{\mathrm{reg}}+\alpha_{\mathrm{reg}}} = \frac{1}{2\sqrt{\alpha_{\mathrm{reg}}}}$ (which also fixes $\norm{S_{w,\alpha_{\mathrm{reg}}}^{-1}} \le (2\sqrt{\alpha_{\mathrm{reg}}})^{-1}$). Solving $e^{-w^2 k_{\mathrm{peak}}^2/2} = \alpha_{\mathrm{reg}}^{1/2}$ yields $w^2 k_{\mathrm{peak}}^2 = \ln(1/\alpha_{\mathrm{reg}}) \implies k_{\mathrm{peak}} = w^{-1}\sqrt{\ln(1/\alpha_{\mathrm{reg}})}$.

\subsection{Parseval Bias--Variance Decomposition and $\mathrm{LG}_{2,1}$ Benchmark --- \josa{18}}
\label{sec:s5_parseval}

\begin{targeteq}{JOSA A Eq.~(18) (p.~7)}
\[
  \mathrm{Bias}^2(\alpha_{\mathrm{reg}}) = \int_0^\infty \left[\frac{\alpha_{\mathrm{reg}}}{e^{-2au}+\alpha_{\mathrm{reg}}}\right]^2 \rho_{2,1}(u)\,\dd u, \qquad
  \mathrm{Var}(\alpha_{\mathrm{reg}}) = \frac{P_\varepsilon}{(w k_{\max})^2}\left[\frac{1}{e^{-w^2 k_{\max}^2}+\alpha_{\mathrm{reg}}} - \frac{1}{1+\alpha_{\mathrm{reg}}}\right].
\]
\end{targeteq}

\paragraph*{Step 1: Splitting total error into deterministic bias and stochastic variance.}
Write $\hat{\psi}_{\alpha_{\mathrm{reg}}} - \psi_{2,1} = (S_{w,\alpha_{\mathrm{reg}}}^{-1}S_w - \mathbb{I})\psi_{2,1} + S_{w,\alpha_{\mathrm{reg}}}^{-1}\varepsilon$. Since $\mathbb{E}[\varepsilon] = 0$, the cross term vanishes upon taking expectation:
\begin{equation}
  \mathrm{MSE}(\alpha_{\mathrm{reg}}) = \norm{(S_{w,\alpha_{\mathrm{reg}}}^{-1}S_w - \mathbb{I})\psi_{2,1}}_{L^2}^2 + \mathbb{E}\!\left[\norm{S_{w,\alpha_{\mathrm{reg}}}^{-1}\varepsilon}_{L^2}^2\right] = \mathrm{Bias}^2 + \mathrm{Var}.
\end{equation}
In Fourier space, $T_{\alpha_{\mathrm{reg}}} - 1 = \frac{g^2}{g^2+\alpha_{\mathrm{reg}}} - 1 = -\frac{\alpha_{\mathrm{reg}}}{g^2+\alpha_{\mathrm{reg}}}$. Using $g^2 = e^{-w^2 k_\perp^2} = e^{-2au}$ (with $u = w_0^2 k_\perp^2/2$ and $a = w^2/w_0^2$) and $\abs{\widetilde{\psi}_{p,l}(\mathbf{k})}^2 \frac{\dd^2\mathbf{k}}{(2\pi)^2} = \rho_{p,l}(u)\,\dd u$ immediately gives the $\mathrm{Bias}^2$ integral in \josa{18}.

\paragraph*{Step 2: Exact closed-form evaluation of $\mathrm{Var}(\alpha_{\mathrm{reg}})$.}
Bandlimited white noise on $\{k_\perp \le k_{\max}\}$ has constant spectral density $\sigma_\varepsilon^2$ satisfying $P_\varepsilon = \int_{k_\perp \le k_{\max}} \sigma_\varepsilon^2 \frac{\dd^2\mathbf{k}}{(2\pi)^2} = \frac{\sigma_\varepsilon^2}{(2\pi)^2}(\pi k_{\max}^2) = \frac{\sigma_\varepsilon^2 k_{\max}^2}{4\pi}$, hence $\sigma_\varepsilon^2 = 4\pi P_\varepsilon / k_{\max}^2$. Integrating $\abs{\widetilde{H}_{w,\alpha_{\mathrm{reg}}}(k_\perp)}^2$:
\begin{equation}
  \mathrm{Var}(\alpha_{\mathrm{reg}}) = \frac{\sigma_\varepsilon^2}{(2\pi)^2}\int_0^{k_{\max}} 2\pi k_\perp\,\frac{e^{-w^2 k_\perp^2}}{\bigl(e^{-w^2 k_\perp^2} + \alpha_{\mathrm{reg}}\bigr)^2}\,\dd k_\perp.
\end{equation}
Substitute $y = e^{-w^2 k_\perp^2}$, so $\dd y = -2w^2 k_\perp e^{-w^2 k_\perp^2}\,\dd k_\perp \implies k_\perp e^{-w^2 k_\perp^2}\,\dd k_\perp = -\frac{\dd y}{2w^2}$, with limits $k_\perp = 0 \mapsto y = 1$ and $k_\perp = k_{\max} \mapsto y = e^{-w^2 k_{\max}^2}$:
\begin{align}
  \mathrm{Var}(\alpha_{\mathrm{reg}}) &= \frac{4\pi P_\varepsilon}{k_{\max}^2}\cdot\frac{2\pi}{4\pi^2}\cdot\frac{1}{2w^2}\int_{e^{-w^2 k_{\max}^2}}^{1}\frac{\dd y}{(y + \alpha_{\mathrm{reg}})^2}
  = \frac{P_\varepsilon}{w^2 k_{\max}^2}\left[-\frac{1}{y + \alpha_{\mathrm{reg}}}\right]_{e^{-w^2 k_{\max}^2}}^{1} \nonumber \\
  &= \frac{P_\varepsilon}{(w k_{\max})^2}\left[\frac{1}{e^{-w^2 k_{\max}^2} + \alpha_{\mathrm{reg}}} - \frac{1}{1 + \alpha_{\mathrm{reg}}}\right],
\end{align}
proving \josa{18} identically!

\begin{pencheck}[Hand Check of Unregularized Blow-Up vs.\ Optimal Tikhonov ($w k_{\max}=3.0, P_\varepsilon=0.01$)]
At $\alpha_{\mathrm{reg}} = 0$, $\mathrm{Bias}^2(0) = 0$ and $\mathrm{MSE}(0) = \mathrm{Var}(0) = \frac{0.01}{3^2}(e^9 - 1) = \frac{0.01}{9}(8103.0839 - 1) = 9.00232$. At $\alpha_{\mathrm{reg}}^{\mathrm{opt}} = 0.049025$: $\mathrm{Var} = \frac{0.01}{9}\bigl(\frac{1}{0.00012341 + 0.049025} - \frac{1}{1.049025}\bigr) = \frac{0.01}{9}(20.3465 - 0.9533) = 0.021548$, $\mathrm{Bias}^2 = 0.013257$, $\mathrm{MSE}_{\min} = 0.034805$, ratio $9.00232/0.034805 = \mathbf{258.65\times}$ ($10\log_{10}258.65 = \mathbf{24.13\,\mathrm{dB}}$).
\end{pencheck}

% ==============================================================================
\section{JOSA A Section 6 \& Appendix A: $\mathrm{LG}_{p,l}$ Expansion, Washout, and Gram Spectrum}
\label{sec:s6}

\subsection{Exact $(p+1)$-Term Closed-Form Expansion $c_{l,p}(R)$ --- \josa{19}--\josa{20} \& \josa{25}--\josa{26}}
\label{sec:s6_clp}

\begin{targeteq}{JOSA A Eqs.~(19), (20), (25), and (26) (pp.~8, 10)}
\[
  c_{l,p}(R) = \mathcal{A}_{p,l}(R)\sum_{m=0}^p (-1)^m \binom{p+\abs{l}}{p-m} \xi^m L_m^{\abs{l}}(-\zeta_R), \qquad
  \xi = \frac{2w^2}{w_0^2+w^2}, \qquad \zeta_R = \frac{R^2 w_0^2}{w^2(w_0^2+w^2)}.
\]
\end{targeteq}

\paragraph*{Step 1: Setting up the radial integral \josa{25}.}
Substitute \josa{1} and $\psi_{p,l}(r,\phi) = N_{p,l}(\sqrt{2}r/w_0)^{\abs{l}} L_p^{\abs{l}}(2r^2/w_0^2) e^{-r^2/w_0^2 + il\phi}$ ($N_{p,l} = \sqrt{2p!/[\pi(p+\abs{l})!]}/w_0$) into $c_{l,p}(R) = \int_0^\infty r\,\dd r\int_0^{2\pi}\dd\phi\,\Phi_{l,R}^* \psi_{p,l}$. The azimuthal integral yields $2\pi$, and combining the Gaussian factors gives $a_0 \equiv \frac{1}{w^2} + \frac{1}{w_0^2} = \frac{w_0^2+w^2}{w^2 w_0^2}$ and $b_R \equiv \frac{2R}{w^2}$, producing \josa{25}:
\begin{equation}
  c_{l,p}(R) = 2\pi C_0 N_{p,l} e^{-R^2/w^2}\left(\frac{\sqrt{2}}{w_0}\right)^{\abs{l}} \int_0^\infty r^{\abs{l}+1} e^{-a_0 r^2} I_{\abs{l}}(b_R r)\,L_p^{\abs{l}}\!\left(\frac{2r^2}{w_0^2}\right)\dd r.
\end{equation}
Check the two key parameter combinations explicitly:
\begin{equation}
  \frac{b_R^2}{4a_0} = \frac{4R^2/w^4}{4(w_0^2+w^2)/(w^2 w_0^2)} = \frac{R^2 w_0^2}{w^2(w_0^2+w^2)} = \zeta_R, \qquad
  -\frac{R^2}{w^2} + \zeta_R = \frac{-R^2(w_0^2+w^2) + R^2 w_0^2}{w^2(w_0^2+w^2)} = -\frac{R^2}{w_0^2+w^2}.
\end{equation}

\paragraph*{Step 2: Explicit polynomial expansion and GR 6.631.1 \josa{26}.}
Expand the associated Laguerre polynomial $L_p^{\abs{l}}(x) = \sum_{m=0}^p \frac{(-1)^m}{m!}\binom{p+\abs{l}}{p-m} x^m$ at $x = 2r^2/w_0^2$:
\begin{equation}
  L_p^{\abs{l}}\!\left(\frac{2r^2}{w_0^2}\right) = \sum_{m=0}^p \frac{(-1)^m}{m!}\binom{p+\abs{l}}{p-m}\left(\frac{2}{w_0^2}\right)^m r^{2m}.
\end{equation}
For each $m$, recall Gradshteyn--Ryzhik formula 6.631.1 ($\operatorname{Re}(a_0)>0$, $\operatorname{Re}(\mu+\nu)>0$):
\begin{equation}
  \int_0^\infty r^{\mu-1} e^{-a_0 r^2} J_\nu(\beta r)\,\dd r = \frac{\beta^\nu\,\Gamma\bigl(\frac{\mu+\nu}{2}\bigr)}{2^{\nu+1} a_0^{(\mu+\nu)/2}\Gamma(\nu+1)}\,{}_1F_1\!\left(\frac{\mu+\nu}{2};\,\nu+1;\,-\frac{\beta^2}{4a_0}\right).
\end{equation}
Continue to $I_\nu(b_R r) = i^{-\nu} J_\nu(i b_R r)$ by setting $\beta = i b_R$: $\beta^\nu = i^\nu b_R^\nu$ cancels $i^{-\nu}$ cleanly, and $-\beta^2/(4a_0) = +b_R^2/(4a_0) = +\zeta_R$. Setting $\nu = \abs{l}$ and $\mu - 1 = \abs{l} + 2m + 1 \implies \mu = \abs{l} + 2m + 2 \implies \frac{\mu+\nu}{2} = \abs{l} + m + 1$ yields \josa{26}:
\begin{equation}
  \int_0^\infty r^{\abs{l}+2m+1} e^{-a_0 r^2} I_{\abs{l}}(b_R r)\,\dd r = \frac{b_R^{\abs{l}}\,\Gamma(\abs{l}+m+1)}{2^{\abs{l}+1} a_0^{\abs{l}+m+1}\Gamma(\abs{l}+1)}\,{}_1F_1(\abs{l}+m+1;\,\abs{l}+1;\,\zeta_R).
\end{equation}

\paragraph*{Step 3: Kummer's first transformation and exact cancellation.}
Apply Kummer's identity ${}_1F_1(\alpha;\,\beta;\,z) = e^z\,{}_1F_1(\beta-\alpha;\,\beta;\,-z)$ with $\alpha = \abs{l}+m+1$, $\beta = \abs{l}+1$, so $\beta - \alpha = -m$:
\begin{equation}
  {}_1F_1(\abs{l}+m+1;\,\abs{l}+1;\,\zeta_R) = e^{\zeta_R}\,{}_1F_1(-m;\,\abs{l}+1;\,-\zeta_R).
\end{equation}
By the hypergeometric representation of the associated Laguerre polynomial, $L_m^{\abs{l}}(y) = \binom{m+\abs{l}}{m}{}_1F_1(-m;\,\abs{l}+1;\,y) = \frac{\Gamma(\abs{l}+m+1)}{m!\,\Gamma(\abs{l}+1)}\,{}_1F_1(-m;\,\abs{l}+1;\,y)$, hence
\begin{equation}
  \frac{\Gamma(\abs{l}+m+1)}{\Gamma(\abs{l}+1)}\,{}_1F_1(-m;\,\abs{l}+1;\,-\zeta_R) = m!\,L_m^{\abs{l}}(-\zeta_R)!
\end{equation}
The factor $m!$ cancels $1/m!$ in the Laguerre expansion, and the $m$-dependent powers simplify to
\begin{equation}
  \left(\frac{2}{w_0^2}\right)^m a_0^{-m} = \left(\frac{2}{w_0^2}\cdot\frac{w^2 w_0^2}{w_0^2+w^2}\right)^m = \left(\frac{2w^2}{w_0^2+w^2}\right)^m = \xi^m.
\end{equation}
Finally collect every $m$-independent factor into $\mathcal{A}_{p,l}(R)$:
\begin{align}
  \mathcal{A}_{p,l}(R) &= 2\pi\left(\frac{1}{w^2}\sqrt{\frac{2}{\pi}}\right)\left(\frac{1}{w_0}\sqrt{\frac{2p!}{\pi(p+\abs{l})!}}\right)\frac{2^{\abs{l}/2}}{w_0^{\abs{l}}}\cdot\frac{(2R/w^2)^{\abs{l}}}{2^{\abs{l}+1}}\cdot\left(\frac{w^2 w_0^2}{w_0^2+w^2}\right)^{\abs{l}+1} e^{-R^2/w^2 + \zeta_R} \nonumber \\
  &= 2\pi\cdot\frac{2}{\pi w^2 w_0}\sqrt{\frac{p!}{(p+\abs{l})!}}\cdot\frac{2^{\abs{l}/2}}{2}\cdot\frac{R^{\abs{l}} w^{2\abs{l}+2} w_0^{2\abs{l}+2}}{w_0^{\abs{l}} w^{2\abs{l}}(w_0^2+w^2)^{\abs{l}+1}}\exp\!\left(-\frac{R^2}{w_0^2+w^2}\right) \nonumber \\
  &= 2^{\abs{l}/2 + 1}\sqrt{\frac{p!}{(p+\abs{l})!}}\,\frac{R^{\abs{l}} w_0^{\abs{l}+1}}{(w_0^2+w^2)^{\abs{l}+1}}\exp\!\left(-\frac{R^2}{w_0^2+w^2}\right),
\end{align}
proving \josa{19}--\josa{20} without a single skipped factor!

\begin{pencheck}[Hand Evaluation of $c_{1,p}(R)$ at $R = 1, w_0 = 1, w = 0.25$ ($p = 0, 1, 2$)]
Here $w_0^2 + w^2 = 1.0625 = 17/16$, $\xi = 2(1/16)/(17/16) = 2/17 = 0.117647$, $\zeta_R = 16/(17/16) = 256/17 = 15.058824$, $e^{-16/17} = 0.390169$.
For $l = 1$: $L_0^1(-\zeta) = 1$; $L_1^1(-\zeta) = 2 + \zeta = 17.058824$; $L_2^1(-\zeta) = 3 + 3\zeta + \zeta^2/2 = 161.560554$.
\begin{itemize}[leftmargin=1.5em,itemsep=1pt,topsep=2pt]
  \item $p=0$: $\mathcal{A}_{0,1} = 2^{3/2}(16/17)^2 e^{-16/17} = 0.977551 \implies c_{1,0}(1) = \mathbf{0.977551}$ (quadrature: $0.977551$).
  \item $p=1$: $\mathcal{A}_{1,1} = \mathcal{A}_{0,1}/\sqrt{2} = 0.691233$; sum $= \binom{2}{1}(1) - \binom{2}{0}\xi L_1^1(-\zeta) = 2 - \frac{2}{17}(17.058824) = -0.0069204 \implies c_{1,1}(1) = 0.691233(-0.0069204) = \mathbf{-0.0047836}$ (quadrature: $-0.0047836$!).
  \item $p=2$: $\mathcal{A}_{2,1} = \mathcal{A}_{0,1}/\sqrt{3} = 0.564390$; sum $= (-1)^0\binom{3}{2}(1) + (-1)^1\binom{3}{1}\xi L_1^1(-\zeta) + (-1)^2\binom{3}{0}\xi^2 L_2^1(-\zeta) = 3 - 3(2.006920) + 2.236132 = -0.784629 \implies c_{1,2}(1) = 0.564390(-0.784629) = \mathbf{-0.442836}$ (exact match to all 6 digits of quadrature: $-0.442836$!).
\end{itemize}
\end{pencheck}

\subsection{Radial Nodal Washout, Exact Rationals, and WKB Arcsine Constant --- \josa{21}--\josa{23}}
\label{sec:s6_washout}

\begin{targeteq}{JOSA A Eqs.~(21), (22), and (23) (p.~8)}
\[
  \mathcal{F}(p,l,w/w_0) = \frac{\bigl(\int_0^\infty e^{-au}\rho_{p,l}(u)\,\dd u\bigr)^2}{\int_0^\infty e^{-2au}\rho_{p,l}(u)\,\dd u}, \quad
  w_{\mathrm{crit}}(0,l) = w_0\left[(1 - \sqrt{1-q_l})^{-1}-1\right]^{1/2}, \quad
  w_{\mathrm{crit}}(p,l) \approx \frac{0.573\,w_0}{\sqrt{2p+\abs{l}+1}}.
\]
\end{targeteq}

\paragraph*{Step 1: Exact rational integration by hand at $w/w_0 = 1$ ($a = 1$), $l = 1$ (Table~2a).}
With $\rho_{p,1}(u) = \frac{1}{p+1} u e^{-u}[L_p^1(u)]^2$ and $\int_0^\infty u^n e^{-\lambda u}\,\dd u = n!/\lambda^{n+1}$:
\begin{itemize}[leftmargin=1.5em,itemsep=1pt,topsep=2pt]
  \item $p = 0$ ($L_0^1 = 1$): $J_1 = \int_0^\infty u e^{-2u}\dd u = \frac{1}{4}$, $J_2 = \int_0^\infty u e^{-3u}\dd u = \frac{1}{9} \implies \mathcal{F}(0,1,1) = \frac{(1/4)^2}{1/9} = \mathbf{\frac{9}{16} = 0.56250}$.
  \item $p = 1$ ($L_1^1 = 2 - u$, $[L_1^1]^2 = 4 - 4u + u^2$, prefactor $1/2$):
    \begin{align*}
      J_1 &= \frac{1}{2}\int_0^\infty (4u - 4u^2 + u^3)e^{-2u}\,\dd u = \frac{1}{2}\left(\frac{4}{4} - \frac{8}{8} + \frac{6}{16}\right) = \frac{3}{16}, \\
      J_2 &= \frac{1}{2}\int_0^\infty (4u - 4u^2 + u^3)e^{-3u}\,\dd u = \frac{1}{2}\left(\frac{4}{9} - \frac{8}{27} + \frac{6}{81}\right) = \frac{1}{2}\cdot\frac{36 - 24 + 6}{81} = \frac{1}{9}, \\
      \mathcal{F}(1,1,1) &= \frac{(3/16)^2}{1/9} = \mathbf{\frac{81}{256} = 0.31641}!
    \end{align*}
\end{itemize}
\paragraph*{Step 2: Exact $p = 0$ algebraic inversion \josa{22}.}
For $p = 0$, $\rho_{0,l}(u) = \frac{1}{\abs{l}!} u^{\abs{l}} e^{-u}$, so $\int_0^\infty e^{-\beta u}\rho_{0,l}(u)\,\dd u = (1+\beta)^{-(\abs{l}+1)}$, yielding $\mathcal{F}(0,l,w/w_0) = \bigl[(1+2a)/(1+a)^2\bigr]^{\abs{l}+1} = \bigl[1 - a^2/(1+a)^2\bigr]^{\abs{l}+1}$. Setting $\mathcal{F} = 0.95$ with $q_l \equiv 0.95^{1/(\abs{l}+1)}$:
\begin{equation}
  1 - \left(\frac{a}{1+a}\right)^2 = q_l \implies \frac{a}{1+a} = \sqrt{1-q_l} \implies \frac{1}{1+a} = 1 - \sqrt{1-q_l} \implies a = \left(1 - \sqrt{1-q_l}\right)^{-1} - 1,
\end{equation}
and taking $w_{\mathrm{crit}} = w_0\sqrt{a}$ yields \josa{22}.

\paragraph*{Step 3: WKB arcsine Laplace transform \josa{23}.}
For $p \gg \abs{l}$, the locally averaged density $\rho_{p,l}(u)$ converges on $u \in (0, 4N_p)$ ($2N_p = 2p + \abs{l} + 1$) to the WKB arcsine density $\rho_{\mathrm{WKB}}(u) = [\pi\sqrt{u(4N_p-u)}]^{-1}$. Substituting $u = 2N_p(1 - \cos\theta)$ ($\theta \in [0,\pi]$, $\dd u = 2N_p\sin\theta\,\dd\theta$, $\sqrt{u(4N_p-u)} = 2N_p\sin\theta$) evaluates the Laplace transform identically:
\begin{equation}
  \int_0^{4N_p} e^{-au}\rho_{\mathrm{WKB}}(u)\,\dd u = \frac{1}{\pi}\int_0^\pi e^{-2N_p a(1-\cos\theta)}\,\dd\theta = e^{-2N_p a} I_0(2N_p a).
\end{equation}
Setting $c \equiv 2N_p a = (2p+\abs{l}+1)(w/w_0)^2$ cancels $e^{-2c}$ identically: $\mathcal{F}_{\mathrm{WKB}}(c) = I_0(c)^2/I_0(2c) = 0.95 \implies c_* = 0.327687 \implies \sqrt{c_*} = \mathbf{0.572440}$, proving \josa{23}.

\subsection{Discrete Radial Sampling and Szeg\H{o} Jacobi-Theta Toeplitz Bound --- \josa{24}}
\label{sec:s6_gram}

\begin{targeteq}{JOSA A Eq.~(24) and Table~2(b) (p.~9)}
\[
  \kappa(T_N) < \kappa(T_\infty) = \frac{\vartheta_3(0,q)}{\vartheta_4(0,q)}, \qquad q = e^{-\Delta R^2/(2w^2)}, \qquad N_{\max} \approx \left\lfloor \frac{R_{\max}-R_{\min}}{2w} \right\rfloor + 1.
\]
\end{targeteq}

For $R_n = R_{\min} + (n-1)\Delta R$ in the annular regime $R_{\min} \gg w, \abs{l}w/2$, Eq.~\josa{6} converges to the symmetric Gaussian Toeplitz matrix $T_{n-m} = q^{(n-m)^2}$ with $q = e^{-\Delta R^2/(2w^2)} \in (0,1)$. For any vector $\mathbf{v} = (v_1,\dots,v_N)^{\mathsf{T}} \in \mathbb{C}^N$, form its trigonometric polynomial $V(\theta) = \sum_{n=1}^N v_n e^{-in\theta}$ ($\frac{1}{2\pi}\int_{-\pi}^\pi \abs{V(\theta)}^2\,\dd\theta = \norm{\mathbf{v}}_2^2$). Writing $q^{(n-m)^2} = \frac{1}{2\pi}\int_{-\pi}^\pi f(\theta) e^{i(n-m)\theta}\,\dd\theta$ with generating symbol
\begin{equation}
  f(\theta) = \sum_{k=-\infty}^\infty q^{k^2} e^{ik\theta} = 1 + 2\sum_{k=1}^\infty q^{k^2}\cos(k\theta) = \vartheta_3\!\left(\frac{\theta}{2}, q\right),
\end{equation}
the Rayleigh quotient of $T_N$ equals
\begin{equation}
  \frac{\mathbf{v}^\dagger T_N \mathbf{v}}{\norm{\mathbf{v}}_2^2} = \frac{\frac{1}{2\pi}\int_{-\pi}^\pi f(\theta)\abs{V(\theta)}^2\,\dd\theta}{\frac{1}{2\pi}\int_{-\pi}^\pi \abs{V(\theta)}^2\,\dd\theta}.
\end{equation}
Since $f(\theta)$ is strictly positive and non-constant (by Jacobi's triple product $\vartheta_3(\theta/2,q) = \prod_{m=1}^\infty (1-q^{2m})(1+2q^{2m-1}\cos\theta + q^{4m-2}) > 0$), every eigenvalue $\lambda_j(T_N)$ satisfies $\min_\theta f(\theta) < \lambda_j(T_N) < \max_\theta f(\theta)$ strictly! Clearly $\max_\theta f(\theta) = f(0) = \vartheta_3(0,q)$ (all $\cos(k\theta) = +1$) and $\min_\theta f(\theta) = f(\pi) = 1 + 2\sum_{k=1}^\infty (-1)^k q^{k^2} = \vartheta_4(0,q)$ (by the triple product where $\cos\theta = -1$ minimizes every factor simultaneously!). Hence $\kappa(T_N) = \lambda_{\max}/\lambda_{\min} < \vartheta_3(0,q)/\vartheta_4(0,q)$ for every $N \ge 1$.

\begin{pencheck}[Hand Arithmetic at $\Delta R = 2.0\,w$ ($q = e^{-2} = 0.135335$)]
Keep just $k = 1, 2$ ($q^4 = 0.0003355$, $q^9 \approx 1.5\times 10^{-8}$):
\begin{align*}
  \vartheta_3(0,e^{-2}) &= 1 + 2(0.135335) + 2(0.0003355) = 1.271341, \\
  \vartheta_4(0,e^{-2}) &= 1 - 2(0.135335) + 2(0.0003355) = 0.729999, \\
  \kappa(T_\infty) &= \frac{1.271341}{0.729999} = \mathbf{1.74157}\quad (\text{bounding }\kappa(\gamma) = 1.6449\text{ at }N=6!).
\end{align*}
Adjacent power isolation is $-10\log_{10}(0.135528^2) = -20\log_{10}(0.135528) = \mathbf{17.36\,\mathrm{dB}}$.
\end{pencheck}

\end{document}
